\section{Introduction}
\label{sec:intro}

Multi-object tracking is dominated by the tracking-by-detection paradigm: a per-frame object detector
proposes boxes, and an association stage links them across time into identity-consistent tracks. DeepSORT
\cite{wojke2017deepsort}, which augments the SORT core \cite{bewley2016sort} (a Kalman filter plus
Hungarian assignment) with a deep appearance embedding, remains a canonical instance of this design and a
common baseline. Its original detector and re-identification (re-ID) encoder are now nearly a decade old, so
the engineering instinct is obvious: swap in a contemporary detector and a modern re-ID backbone behind the
unchanged SORT core. We do exactly this, integrating pluggable detectors (YOLOv8m \cite{jocher2023yolov8})
and re-ID models (OSNet \cite{zhou2019osnet}; a generic ImageNet timm backbone \cite{wightman2019timm})
through a single \texttt{Detection(tlwh, confidence, feature)} contract. On six MOT-Challenge training videos
\cite{leal2015mot15,milan2016mot16}, mean HOTA \cite{luiten2021hota,luiten2020trackeval} rises from
\baselineHOTA{} to \modernHOTA{} ($+\deltaHOTA$) at \fpsRealtime{} FPS. This headline result is also
unsurprising and, on its own, unpublishable: a stronger detector and a stronger encoder make a better
tracker.

The interesting questions are the ones the headline number hides. \emph{Where} does the gain actually come
from -- the better detector, or the better appearance model? And does a stronger detector always help: can
\emph{more} detections, obtained by raising recall, ever \emph{hurt} tracking quality? We find that the
answers are not the intuitive ones. The modernization gain is overwhelmingly detection-driven, and on dense,
low-resolution clips pushing the detector toward higher recall drives HOTA monotonically \emph{down}. The
mechanism is concrete and reproducible: the modern live pipeline silently drops the original DeepSORT gating
contract (a $0.3$ confidence threshold plus non-maximum suppression) that once filtered the precomputed
detections, so ungated false positives instantiate spurious tracks. This leads to our thesis:
\textbf{in tracking-by-detection, per-clip detector-confidence (precision) gating is a larger and more
reliable lever on HOTA than global recall maximization}.

We make four contributions.
\begin{enumerate}
\item A detection-vs-association decomposition of the modernization gain, showing it is
\textbf{detection-dominated}: on MOT16 the detection term improves by $\detaGain$ HOTA against only
$\assaGain$ for association, and under ground-truth boxes a non-re-ID-trained ImageNet backbone reaches
\timmGT{} HOTA versus \osnetGT{} for a dedicated re-ID model, so association is nearly detector-bound once
boxes are fixed.
\item A counterintuitive \textbf{``more detections hurt''} regime on dense low-resolution clips: raising the
TUD-Campus input resolution lowers HOTA monotonically
($\imgszLow\!\rightarrow\!\imgszMid\!\rightarrow\!\imgszHigh$), and the sign of recall's effect is predicted
by per-clip detector precision across all \nVideos{} videos.
\item A gating-restoration ablation that attributes the regression to the dropped confidence and NMS gate
and recovers the loss.
\item A practical \textbf{per-clip precision-gating guideline}: tuning the confidence threshold per clip
lifts the worst sequence from a $\tudUntuned$ deficit to $+\tudGain$ HOTA over baseline.
\end{enumerate}

We position this finding against ByteTrack \cite{zhang2022bytetrack} and other modern DeepSORT descendants
\cite{du2023strongsort,cao2023ocsort,aharon2022botsort} in \cref{sec:related}, detail the system and
protocol in \cref{sec:method}, report the audit in \cref{sec:results}, and discuss limitations and
reproducibility lessons in \cref{sec:limitations} before concluding.
